\documentclass[12pt]{article}

% Core math
\usepackage{amsmath, amssymb, amsthm}

% Diagrams
\usepackage{tikz-cd}
\usepackage{tikz}
\usetikzlibrary{decorations.markings, arrows.meta, positioning}

% References
\usepackage{hyperref}
\usepackage{cleveref}
\usepackage{url}

% Graphics
\usepackage{graphicx}

% Page layout
\usepackage[margin=1in]{geometry}

% GrokRxiv sidebar (uses LaTeX 2020+ kernel shipout hook in lieu of
% the legacy `everypage` package).
\usepackage{xcolor}

% Listings for code blocks
\usepackage{listings}
\lstset{
  basicstyle=\ttfamily\footnotesize,
  columns=fullflexible,
  breaklines=true,
  showstringspaces=false,
  frame=single,
  framerule=0.4pt,
  rulecolor=\color{black!40},
  xleftmargin=0pt,
  xrightmargin=0pt
}

% Theorem environments
\newtheorem{theorem}{Theorem}[section]
\newtheorem{proposition}[theorem]{Proposition}
\newtheorem{lemma}[theorem]{Lemma}
\newtheorem{corollary}[theorem]{Corollary}
\theoremstyle{definition}
\newtheorem{definition}[theorem]{Definition}
\newtheorem{example}[theorem]{Example}
\theoremstyle{remark}
\newtheorem{remark}[theorem]{Remark}

% Common HoTT macros
\newcommand{\Type}{\mathsf{Type}}
\newcommand{\U}{\mathcal{U}}
\newcommand{\N}{\mathbb{N}}
\newcommand{\R}{\mathbb{R}}
\newcommand{\Z}{\mathbb{Z}}
\newcommand{\Q}{\mathbb{Q}}
\newcommand{\Rc}{\mathbb{R}_{c}}
\newcommand{\Rd}{\mathbb{R}_{d}}
\newcommand{\Id}{\mathsf{Id}}
\newcommand{\refl}{\mathsf{refl}}
\newcommand{\zero}{\mathsf{zero}}
\newcommand{\suc}{\mathsf{succ}}
\newcommand{\transport}{\mathsf{transport}}
\newcommand{\ua}{\mathsf{ua}}
\newcommand{\idtoeqv}{\mathsf{idtoeqv}}
\newcommand{\eqv}{\simeq}
\newcommand{\fst}{\mathsf{fst}}
\newcommand{\snd}{\mathsf{snd}}
\newcommand{\inl}{\mathsf{inl}}
\newcommand{\inr}{\mathsf{inr}}
\newcommand{\Path}{\mathsf{Path}}
\newcommand{\PathP}{\mathsf{PathP}}
\newcommand{\rat}{\mathsf{rat}}
% NOTE: we deliberately avoid redefining the standard \lim primitive.
\newcommand{\hlim}{\mathsf{lim}}
\newcommand{\II}{\mathbb{I}}
\newcommand{\hcomp}{\mathsf{hcomp}}
\newcommand{\transp}{\mathsf{transp}}
\newcommand{\Glue}{\mathsf{Glue}}
\newcommand{\isContr}{\mathsf{isContr}}
\newcommand{\isProp}{\mathsf{isProp}}
\newcommand{\isSet}{\mathsf{isSet}}
\newcommand{\isEquiv}{\mathsf{isEquiv}}
\newcommand{\Cauchy}{\mathsf{Cauchy}}
% Closeness relation: '~'. Always written with an explicit subscript at
% the call site, e.g. $u \close_{\varepsilon} v$, since the precision is
% an essential part of the relation and we want it visible.
\newcommand{\close}{\mathrel{\sim}}

% GrokRxiv sidebar configuration. We use the LaTeX 2020+ kernel
% shipout hook ('shipout/background') in preference to the legacy
% `everypage` package, which emits deprecation warnings in modern
% LaTeX distributions.
\definecolor{grokgray}{RGB}{110,110,110}
\AddToHook{shipout/background}{%
  \ifnum\value{page}=1
    \begin{tikzpicture}[remember picture, overlay]
      \node[
        rotate=90,
        anchor=south,
        font=\Large\sffamily\bfseries\color{grokgray},
        inner sep=0pt
      ] at ([xshift=38pt, yshift=0.52\paperheight]current page.south west)
      {GrokRxiv:2026.05.cubical-hiit-cauchy\quad
       [\,math.LO\,]\quad
       04 May 2026};
    \end{tikzpicture}
  \fi
}

\title{Cubical Higher Inductive--Inductive Types\\
       and the Cauchy Reals\\
\large A Cubical Agda Completion of the Book HoTT Construction}
\author{YonedaAI Research \\
\textit{Univalent Correspondence Project} \\
\textit{Magneton Foundational Studies} \\
\texttt{mlong168@gmail.com}}
\date{4 May 2026}

\begin{document}
\maketitle

\begin{abstract}
The Cauchy reals $\Rc$ admit a higher inductive--inductive presentation in
Book HoTT (HoTT Book \S 11.3, Booij 2020), and this presentation underwrites
the unique-existence definitions of $\pi$ and $e$ used throughout the
\emph{Univalent Correspondence} series (Paper~V \S 8). However, the
construction is given in Book HoTT, where path constructors are
\emph{postulated} together with their associated $\beta$-rules and where
function extensionality is derived from univalence. A computational analogue
in Cubical Agda---in which paths are functions $\II \to A$, univalence
reduces, and higher inductive types compute---has been
\emph{in progress but incomplete}: the existing Cubical Agda standard
library exposes set quotients and several truncations, but not a clean
HIIT presentation of $\Rc$ with the closeness-relation constructor.

This paper completes the cubical version. We give a Cubical
Agda specification of the simultaneous higher inductive--inductive type
$(\Rc, \close)$, where $\close : \Q^{>0} \to \Rc \to \Rc \to \Type$ is the
$\varepsilon$-closeness predicate of Booij. The four point/path/square
constructors are expressed as primitive cubical operations: the limit
constructor uses $\PathP$ rather than \texttt{Path}, the closeness path
constructor uses an interval-typed family, and the propositional
truncation of $\close$ is replaced by a $\PathP$-valued square
constructor that we prove satisfies the same universal property as the
truncated Book HoTT version. We verify that the cubical $\Rc$ is an h-set,
that it carries the structure of an Archimedean ordered field, and that it
is equivalent (via a cubical \texttt{Glue} type) to the Dedekind reals
$\Rd$ when the latter is restricted to located Dedekind cuts. We also
extract, via Cubical Agda's reduction discipline, an executable map
$\mathsf{run} : \mathbb{N} \to \Q$ approximating $\sqrt{2}$, $\pi$, and $e$;
the rational approximants we obtain by reducing the constructively-defined
unique-existence centres are checked numerically against Haskell prototypes
(Section~7). We close with the open problems that remain: judgemental
$\eta$-rules for the relator constructor, full coherence of the IIT
(inductive--inductive) elimination principle in the presence of Glue
types, and integration with the \texttt{Cubical.HITs} hierarchy
(\texttt{Cubical.HITs.CauchyReals} as a candidate library entry).
\end{abstract}

\tableofcontents

\section{Introduction}
\label{sec:intro}

The constructive real numbers occupy a privileged position in the
foundations of analysis. In a setting without the law of excluded middle,
no canonical definition of $\R$ exists; instead, several non-equivalent
candidates compete: \emph{Cauchy reals} $\Rc$, defined as a quotient of
Cauchy sequences of rationals; \emph{Dedekind reals} $\Rd$, defined as
located cuts; and \emph{streams} or \emph{redundant binary signed-digit
representations}, definable as final coalgebras (Paper~III, \S 5). In
classical mathematics these collapse to a single object, but constructively
they differ by countable choice, fan-theorem-like principles, or
modulus-of-convergence considerations~\cite{BridgesRichman}. The
\emph{Cauchy} presentation is the most useful for executable analysis: it
admits a direct definition of limits, supports the unique-existence
definitions of transcendentals via Picard-style fixed points, and yields
an Archimedean ordered field on the nose.

In Homotopy Type Theory (HoTT)~\cite{HoTTBook}, an additional dimension is
available: the \emph{path} structure of identity types. \S 11.3 of the
HoTT Book, refined by Booij~\cite{Booij2020}, defines $\Rc$ not as a
\emph{quotient} but as a \emph{higher inductive--inductive type} (HIIT):
the carrier $\Rc$ and a parameterised closeness relation
$\close_{\varepsilon} : \Rc \to \Rc \to \Type$ are introduced
\emph{simultaneously}, with constructors that include both points and
\emph{path constructors} identifying close-enough representatives. The
result is an h-set on which the limit operation
$\hlim : (\Q^{>0} \to \Rc) \to \Rc$ exists \emph{by construction}, with no
choice principle invoked.

This Book HoTT construction is the foundation of Paper~V's treatment of
$\pi$ and $e$ (Definitions 8.1, 8.2, here recalled in Section~\ref{sec:hott-recap}),
which expresses each as the centre of a contractible type of solutions to
its characteristic structural condition. As long as $\Rc$ is a
\emph{type}---and not merely a proposition---these definitions are
mathematically meaningful even before one has computational content for
the rationals approximating $\pi$ to a given precision.

\paragraph{The gap.}
The synthesis of the prior series~\cite[\S 8, item~6]{SynthesisHoTT} flagged the following gap:

\begin{quote}\itshape
The Cauchy reals construction of Paper~V is given in Book HoTT; a clean
cubical version (in Cubical Agda) is in progress but incomplete.
\end{quote}

The gap matters for three reasons. First, Book HoTT is a postulational
extension of MLTT: path constructors are added with their $\beta$-rules
postulated, breaking canonicity in the absence of a model with a
computational interpretation. Cubical type theory (CCHM,
\cite{CCHM}) restores canonicity by introducing the
\emph{interval} $\II$ as a primitive judgemental object with a de~Morgan
algebra structure and replacing path postulates with operations on
$\II$-functions. Second, in Cubical Agda~\cite{CubicalAgda}, univalence is
\emph{not} an axiom but a \emph{theorem}: $\ua$ is a defined term whose
$\beta$-rule reduces. The same machinery suggests that $\Rc$ should be
expressible as an honest cubical HIIT, with all of its eliminators
computing. Third, downstream applications (transcendentals, complex
analysis, $\zeta$-as-HoTT-native-statement; cf.\ \cite[\S 7.3]{SynthesisHoTT}) need a
\emph{computational} $\Rc$, not a postulated one.

\paragraph{What we do.}
We give a Cubical Agda specification of $(\Rc,\close)$ as a HIIT with four
constructors: \texttt{rat}, \texttt{lim}, \texttt{eq} (closeness-induced
path), and a square constructor expressing propositional truncation of
$\close$. The novelty over the Book HoTT presentation is that
\emph{every} constructor is given as either a $\Type$, a $\Path$, or a
$\PathP$ (a path over a path); we never invoke set-truncation or
propositional-truncation as black boxes. The closeness predicate is
defined simultaneously by an inductive-inductive scheme whose well-formedness
in Cubical Agda we justify by reduction to the inductive-inductive
schema of Altenkirch--Kaposi~\cite{AltenkirchKaposi}.

We then prove four results, \emph{cubically}:

\begin{enumerate}
\item[(R1)] \textbf{$\Rc$ is an h-set.} (Theorem~\ref{thm:Rc-isSet})
\item[(R2)] \textbf{$\Rc$ has the universal property of the Cauchy
completion of $\Q$.} (Theorem~\ref{thm:Rc-univ})
\item[(R3)] \textbf{There is a cubical equivalence
$\Rc \eqv \Rd^{\text{loc}}$, where $\Rd^{\text{loc}}$ is the type of
located Dedekind cuts.} (Theorem~\ref{thm:dedekind})
\item[(R4)] \textbf{$\Rc$ extracts: the rational approximants to
$\sqrt{2}$, $\pi$, and $e$ are computed by reduction of normal forms in
Cubical Agda.} (Section~\ref{sec:extraction})
\end{enumerate}

The remainder of the paper makes (R1)--(R4) precise, gives the Cubical
Agda code for the central definitions and theorems, and identifies the
remaining gaps---chiefly the absence of a fully judgemental $\eta$-rule
for the closeness constructor and the lack of an integrated
\texttt{Cubical.HITs.CauchyReals} module in the standard library
\cite{CubicalAgdaLib}.

\paragraph{Outline.}
Section~\ref{sec:hott-recap} recalls the Book HoTT construction of $\Rc$
from~\cite{HoTTBook,Booij2020}.
Section~\ref{sec:cchm} reviews CCHM cubical type theory: the interval, Kan
operations, Glue types, computational univalence.
Section~\ref{sec:cubical-hiits} discusses cubical HIITs in general,
focusing on what changes when one moves from Book HoTT path-constructors
to Cubical Agda \texttt{Path}/\texttt{PathP} constructors.
Section~\ref{sec:cubical-cauchy} gives the Cubical Agda HIIT of $\Rc$.
Section~\ref{sec:dedekind} sketches the equivalence with located Dedekind
cuts.
Section~\ref{sec:extraction} discusses computational content and
extraction of rational approximants.
Section~\ref{sec:open} lists open problems and the path forward.

\paragraph{Audience.}
We assume familiarity with HoTT at the level of~\cite{HoTTBook}, Chapters
1--3 and~6, and basic constructive analysis at the level of
\cite{BridgesRichman}. Cubical type theory at the level of~\cite{CCHM} is
helpful but is reviewed here. Familiarity with Cubical Agda
\cite{CubicalAgda} is helpful for reading the code blocks but is not
strictly required.

\section{Mathematical Framework and Notation}
\label{sec:framework}

This section fixes notation and recalls the type-theoretic vocabulary
used throughout. Readers familiar with the HoTT Book may safely skip
to Section~\ref{sec:hott-recap}.

\subsection{Universes and Truncation Levels}

We work in a tower of universes
$\U_0 : \U_1 : \U_2 : \dots$, with the Tarski-style cumulativity
$\U_n \subseteq \U_{n+1}$. We write $\Type$ for the
generic universe; the index is left implicit unless required.

For $X : \Type$, the predicates
$\isContr(X), \isProp(X), \isSet(X)$ are defined as in
\cite[\S 3]{HoTTBook}:
\begin{align*}
\isContr(X) &:\equiv \Sigma_{x : X}\, \Pi_{y : X}\, x = y, \\
\isProp(X)  &:\equiv \Pi_{x, y : X}\, x = y, \\
\isSet(X)   &:\equiv \Pi_{x, y : X}\, \isProp(x = y).
\end{align*}
A type $X$ with $\isSet(X)$ is called an \emph{h-set} (homotopy set). The
type $\Rc$ to be defined will be an h-set, but its definition crucially
uses the path structure between paths: it is not just a quotient.

\subsection{Identity Types and Path Types}

In Book HoTT, the identity type $a =_A b$ is generated by reflexivity
and eliminated by the J-rule. In Cubical Agda, the type
$\Path\,A\,a\,b$ replaces this and is judgmentally
$\II \to A$ with $a, b$ at the endpoints. The two are
equivalent (\cite[\S 6]{CCHM}), but the cubical formulation makes
function extensionality, univalence, and HIT $\beta$-rules computable.

\begin{definition}[Equivalence]
\label{def:equiv}
For $f : A \to B$, $\isEquiv(f) :\equiv \Pi_{b : B}\, \isContr(\Sigma_{a : A}\, f\,a = b)$.
The type of equivalences is $A \eqv B :\equiv \Sigma_{f : A \to B}\, \isEquiv(f)$.
\end{definition}

\subsection{Higher Inductive Types and HIITs}

A \emph{higher inductive type} (HIT) is presented by a list of point
constructors and path constructors (paths in the resulting type), with
an eliminator that respects all of these. The canonical example is the
circle $S^1$, with $\mathsf{base} : S^1$ and
$\mathsf{loop} : \mathsf{base} = \mathsf{base}$.

A \emph{higher inductive--inductive type} (HIIT) introduces both a
type $A$ and a family $B$ over $A$ (or over $A \times A$, etc.) with
constructors that may reference inhabitants of either. The
formal schema for HIITs in Book HoTT is given by Altenkirch--Kaposi
\cite{AltenkirchKaposi} via quotient inductive--inductive types
(QIITs); the cubical version is treated by~\cite{CCHM} for QITs and
extended in~\cite{CubicalAgda} to handle the HIIT case.

\subsection{Conventions}

\begin{itemize}
\item $\Pi_{x : A}\, B(x)$ is the dependent product (function) type,
also written $(x : A) \to B(x)$ in Agda syntax.
\item $\Sigma_{x : A}\, B(x)$ is the dependent sum (pair) type.
\item We write $f \circ g$ for function composition,
$f^{-1}$ for the inverse of an equivalence.
\item We use the convention that bound variables in
$\Pi$ and $\Sigma$ types may be introduced silently when no ambiguity
arises (so we write $\Pi q.\, P\,q$ for $\Pi_{q : T}\, P\,q$ when $T$
is clear from context).
\end{itemize}

\section{The Book HoTT Cauchy Reals: Recap of Paper V}
\label{sec:hott-recap}

We follow~\cite[\S 11.3]{HoTTBook}, reorganising slightly to align with
Booij's PhD thesis~\cite{Booij2020}, which is the most detailed source
for this construction.

\subsection{Rationals and Closeness}

We assume the type $\Q$ of rational numbers is given as a set quotient
$(\Z \times \Z_{>0}) /\!\!\sim$, where $\sim$ is the cross-multiplication
equivalence. We write $\Q^{>0}$ for the type of strictly positive
rationals, equivalently $\Sigma_{q : \Q}\, q > 0$. The standard
operations $+, -, \cdot, |\cdot|$ are inherited from $\Q$; we treat them
as judgmentally well-defined.

\begin{definition}[$\varepsilon$-closeness on $\Q$]
\label{def:close-Q}
For $\varepsilon : \Q^{>0}$ and $p, q : \Q$, define
$p \close_{\varepsilon} q :\equiv |p - q| < \varepsilon$.
\end{definition}

\subsection{The HIIT Presentation}

A brief orientation. The Cauchy condition we use,
$x_{\delta} \close_{\delta + \varepsilon} x_{\varepsilon}$ for all
$\delta, \varepsilon : \Q^{>0}$, is the
\emph{constructive} formulation: a sequence carries with it an
\emph{explicit modulus of convergence}, given by the index
$\varepsilon$ itself. Classically, one writes
``$\forall \varepsilon\,\exists N\,\forall n \geq N\,|x_n - x_m| <
\varepsilon$,'' but this requires both choice (to extract $N$ from the
existence claim) and a $\N$-indexed sequence. We avoid both: the
sequence is indexed by $\Q^{>0}$ (so the rate of convergence is
encoded in the indexing itself), and the Cauchy condition is
$\Pi$-typed (so the modulus is an internal datum, not extracted via
choice). This is the formulation adopted by Booij
\cite{Booij2020} and by the HoTT Book
\cite[\S 11.3]{HoTTBook}.

\begin{definition}[Cauchy reals as a HIIT, Booij]
\label{def:Rc-book}
The pair $(\Rc, \close)$ is the simultaneous higher
inductive--inductive type with the following constructors.

\textbf{Carrier constructors of $\Rc$:}
\begin{itemize}
\item $\rat : \Q \to \Rc$.
\item $\hlim : \Pi_{x : \Q^{>0} \to \Rc}\,
       \mathsf{isCauchyApprox}(x) \to \Rc$, where
\[
\mathsf{isCauchyApprox}(x) :\equiv
\Pi_{\delta, \varepsilon : \Q^{>0}}\,
   x_{\delta} \close_{\delta + \varepsilon} x_{\varepsilon}.
\]
\item $\mathsf{eq}_{\Rc} : \Pi_{u, v : \Rc}\,
        (\Pi_{\varepsilon : \Q^{>0}}\, u \close_{\varepsilon} v) \to u = v$.
\end{itemize}

\textbf{Constructors of $\close_{\varepsilon}$ (one per pair-shape).}
A small but crucial detail: the constructors below take subtractions
of strictly positive rationals as $\Q^{>0}$-precisions, so each is
implicitly conditioned on the relevant strict inequality
($\varepsilon > \delta$, $\varepsilon > \delta + \eta$). We carry these
preconditions along throughout the paper; equivalently, one may
reparameterise so that the gap is given as an additional
$\Q^{>0}$ argument.
\begin{itemize}
\item \textbf{$\mathsf{rat\text{-}rat}$.}
\[
\Pi_{q, r : \Q}\,\Pi_{\varepsilon : \Q^{>0}}\,
   |q - r| < \varepsilon \;\longrightarrow\;
   \rat\,q \close_{\varepsilon} \rat\,r.
\]
\item \textbf{$\mathsf{rat\text{-}lim}$.}
\[
\begin{aligned}
&\Pi_{q : \Q}\,\Pi_{y, \mathsf{cy}}\,
\Pi_{\varepsilon, \delta : \Q^{>0}}\,
\Pi_{h_{>} : \delta < \varepsilon}, \\
&\quad \rat\,q \close_{\varepsilon - \delta} y_{\delta} \;\longrightarrow\;
   \rat\,q \close_{\varepsilon} \hlim(y, \mathsf{cy}).
\end{aligned}
\]
\item \textbf{$\mathsf{lim\text{-}rat}$.} Symmetric, with precondition
$\delta < \varepsilon$.
\item \textbf{$\mathsf{lim\text{-}lim}$.}
\[
\begin{aligned}
&\Pi_{x, \mathsf{cx}, y, \mathsf{cy}}\,
\Pi_{\varepsilon, \delta, \eta : \Q^{>0}}\,
\Pi_{h_{>} : \delta + \eta < \varepsilon}, \\
&\quad x_{\delta} \close_{\varepsilon - \delta - \eta} y_{\eta}
   \;\longrightarrow\;
   \hlim(x,\mathsf{cx}) \close_{\varepsilon} \hlim(y,\mathsf{cy}).
\end{aligned}
\]
\end{itemize}
The preconditions are essential: $\Q^{>0}$ is not closed under
subtraction, and $\varepsilon - \delta$ is well-defined in $\Q^{>0}$
exactly when $\delta < \varepsilon$. Booij's thesis
\cite[Definition~6.2.1]{Booij2020} carries these inequalities as part
of the constructor signatures; we adopt the same convention. An
alternative formulation, equivalent up to coercion, takes the
\emph{positive gap} $\theta : \Q^{>0}$ as an explicit argument and
writes $\varepsilon = \delta + \theta$ in place of $\varepsilon - \delta$;
this avoids the precondition entirely.

\textbf{Path constructor of $\close$:}
\begin{itemize}
\item For each $\varepsilon, u, v$, the type $u \close_{\varepsilon} v$ is
a proposition: any two inhabitants are equal.
\end{itemize}

\textbf{Set-truncation:}
\begin{itemize}
\item $\Rc$ is an h-set: for any $u, v : \Rc$ and $p, q : u = v$, $p = q$.
\end{itemize}
\end{definition}

The dependent eliminator for $(\Rc, \close)$ provides simultaneous
induction principles for any motive $A : \Rc \to \Type$ and
$B : \Pi_{u, v : \Rc}\, A(u) \to A(v) \to \Q^{>0} \to \Type$
respecting all the constructors and path constructors. For a full
statement, see~\cite[Lemma 11.3.10]{HoTTBook}.

\subsection{Universal Property}

\begin{theorem}[Universal property of $\Rc$,~\cite{HoTTBook} Thm.~11.3.5]
\label{thm:UP-book}
For any complete metric space $Y$ over $\Q$ that is an h-set, every
$\Q$-uniformly continuous map $f : \Q \to Y$ extends uniquely to
$\bar f : \Rc \to Y$.
\end{theorem}

This is the Cauchy-completion universal property and the core technical
result needed for the unique-existence definitions of $\pi$ and $e$.

\subsection{Why a HIIT and Not a Quotient}

In \emph{set-theoretic} foundations one defines $\Rc$ as a quotient of
Cauchy sequences of rationals modulo the equivalence
$x \sim y \iff \forall \varepsilon\, \exists N\, \forall n \geq N\,
  |x_n - y_n| < \varepsilon$. In type theory \emph{without} countable
choice (e.g.\ propositions-as-types MLTT), this quotient is \emph{not}
Cauchy-complete: lifting a Cauchy sequence of equivalence classes to a
Cauchy sequence of representatives requires the choice principle. The HIIT
construction of Definition~\ref{def:Rc-book} avoids this problem by
introducing $\hlim$ as a \emph{primitive constructor} and identifying
representatives via the path constructor $\mathsf{eq}_{\Rc}$.

\subsection{\texorpdfstring{$\pi$ and $e$}{pi and e} as Unique Existence Centres}

\begin{definition}[$\pi$, after~\cite{HoTTBook} 11.3 and Paper V Def.~8.1]
\label{def:pi}
$\pi : \Rc$ is the centre of the contractible type
\[
P_{\pi} :\equiv
\Sigma_{p : \Rc}\,
 (\sin p =_{\Rc} 0) \times (p > 0) \times
 (\Pi_{x : \Rc}\, 0 < x < p \to \sin x > 0).
\]
Here $\sin : \Rc \to \Rc$ is the unique solution of
$\sin'' + \sin = 0,\, \sin\,0 = 0,\, \sin'\,0 = 1$, again presented as a
unique-existence centre.
\end{definition}

\begin{definition}[$e$, after Paper V Def.~8.2]
\label{def:e}
$e :\equiv \exp\,1$, where $\exp : \Rc \to \Rc$ is the centre of
\[
P_{\exp} :\equiv
\Sigma_{f : \Rc \to \Rc}\,
 (f' = f) \times (f\,0 = 1).
\]
\end{definition}

The crucial property of these definitions is that, although $\pi$ and $e$
are introduced as \emph{centres of contractible types}, the centres are
\emph{terms of type $\Rc$}; in Book HoTT, this term is well-typed but
\emph{has no canonical form}. To extract approximating rationals one must
either move outside Book HoTT (an external interpretation) or move to a
type theory where $\Rc$ has computational content. The latter route is
the one we pursue here.

\subsection{What Cubical Has That Book HoTT Lacks}

\begin{itemize}
\item \textbf{Computational univalence.} In Book HoTT, $\ua$ is an axiom;
$\transport$ along $\ua\,(e : A \eqv B)$ does \emph{not} reduce to
$e$ definitionally. Cubical Agda makes this reduction judgmental
(\cite[Thm.~7.2]{Paper05HoTT}).
\item \textbf{HIT canonicity.} The $\beta$-rules of HIT path constructors
in Book HoTT hold only propositionally. In Cubical Agda, the rules hold
\emph{definitionally} for points and \emph{up to higher coherence} for
paths; this is the source of the canonicity theorem of Huber~\cite{Huber}.
\item \textbf{Function extensionality.} A theorem in cubical type theory,
not an axiom: \texttt{funExt} is the application of the path operation to
$i \mapsto \lambda x.\, p\,x\,i$.
\end{itemize}

\section{CCHM Cubical Type Theory}
\label{sec:cchm}

We summarise the features of CCHM~\cite{CCHM} that we will use; this
section is meant as orientation, not a self-contained development.

\subsection{The Interval}

The CCHM type theory adds a primitive object $\II$, the \emph{interval},
together with constants $0, 1 : \II$ and de~Morgan algebra operations
$\sqcap, \sqcup : \II \to \II \to \II$ and $\neg : \II \to \II$
satisfying:
\[
i \sqcap 0 = 0,\quad i \sqcup 1 = 1,\quad \neg(\neg i) = i,\quad
\neg(i \sqcap j) = \neg i \sqcup \neg j, \ldots
\]
The interval is \emph{not} a type in the universe; rather, it is a
sort with its own variable bindings.

\subsection{Path Types}

For $A : \Type$ and $a, b : A$, the path type $\Path_A\,a\,b$ is the type
of functions $p : \II \to A$ with $p\,0 \equiv a$ and $p\,1 \equiv b$
\emph{judgmentally}. Identity is recovered by abbreviation.

\begin{definition}[$\Path$ vs $\PathP$]
\label{def:path-pathp}
$\PathP\,A\,a\,b$, where $A : \II \to \Type$, is the type of dependent
paths: $p : (i : \II) \to A\,i$ with $p\,0 \equiv a : A\,0$ and
$p\,1 \equiv b : A\,1$. We have $\Path_X = \PathP\,(\lambda i.\,X)$.
\end{definition}

\subsection{\texorpdfstring{Kan Operations: $\hcomp$ and $\transp$}{Kan Operations: hcomp and transp}}

The two primitive Kan operations are:
\begin{itemize}
\item $\hcomp$: \emph{homogeneous composition}. Given a partial element
of $A$ defined on a set of faces $\varphi$ of an open box, together with a
``floor'' agreeing on the intersection, $\hcomp$ produces an element of
$A$ giving the lid of the box.
\item $\transp$: \emph{transport along a path of types}. Given
$A : \II \to \Type$ and $a : A\,0$, $\transp\,A\,a : A\,1$.
\end{itemize}

These two operations together implement what classical homotopy theory
calls \emph{Kan filling}: every horn has a lifting to a simplex.
\subsection{Glue Types and Univalence}

\begin{definition}[Glue, sketched]
\label{def:glue}
Given $A : \Type$ and a partial type $T$ defined on a face $\varphi$
together with a partial equivalence $f : T \eqv A$ on $\varphi$, the
type $\Glue\,A\,[\varphi \mapsto (T, f)]$ extends $A$ outside $\varphi$
and equals $T$ on $\varphi$. There is a definable map
$\mathsf{unglue} : \Glue\,A\,[\dots] \to A$ that on $\varphi$ agrees
with $f$.
\end{definition}

\begin{theorem}[Cubical univalence,~\cite{CCHM} Thm.~7.2]
\label{thm:cubical-ua}
For $A, B : \Type$ and $e : A \eqv B$, define
\[
\ua\,e :\equiv \lambda i.\, \Glue\,B\,[(i = 0) \mapsto (A, e),\,
                                   (i = 1) \mapsto (B, \mathsf{idEqv})].
\]
Then $\transport^{\ua\,e}\,a \equiv e\,a$ definitionally, and
$\idtoeqv \circ \ua \equiv \mathsf{id}$.
\end{theorem}

This is the \emph{computational} content of univalence: the inverse of
$\idtoeqv$ is given by the $\Glue$ construction, and the $\beta$-rule
\emph{reduces}.

\subsection{Coercion and Filling}

A subtlety we shall need: the operation $\transp$ in CCHM is governed
by a face condition $\varphi$ that controls when transport is the
identity. A typical signature is
\[
\transp : (A : \II \to \Type)\,(\varphi : \mathbb{F})\,(a_0 : A\,0) \to A\,1,
\]
with the rule that if $\varphi$ is true then $A$ is constant on $\II$
and $\transp$ is the identity. This is what makes $\transp$ behave
like Kan filling rather than being a black box. For our purposes, the
relevant consequence is that transport along the constant path is the
identity \emph{judgmentally}; this is critical for the proofs of
Theorems~\ref{thm:Rc-isSet} and \ref{thm:Rc-univ}, where many path
constructors are filled by reflexivity.

\subsection{The de Morgan Structure on the Interval}

The interval $\II$ in CCHM has a de Morgan algebra structure:
$\sqcap, \sqcup : \II \times \II \to \II$ satisfy associativity,
commutativity, idempotence, distributivity, and de Morgan's laws via
$\neg$. This is \emph{strictly weaker} than the Boolean algebra
structure that one would get if the interval had complementation:
$i \sqcup \neg i \neq 1$ in general. The de Morgan structure is
sufficient for cubical type theory because all we need is the closure
of the Kan operations under face composition; we do not need a unique
``opposite'' to each face.

\begin{remark}[Cartesian vs CCHM cubical]
\label{rem:cartesian}
An alternative cubical type theory, the \emph{cartesian} cubical
theory of Angiuli--Brunerie--Coquand--Harper--Hou--Licata
\cite{ABCHHL}, drops the de Morgan operations and uses only
$\sqcap, \sqcup$ as primitive. The two theories are equivalent at the
level of definitional equality on closed terms, but differ in how
diagonal squares (paths $i \mapsto p\,i\,i$) are formed. Cubical Agda
uses the CCHM presentation; we follow that throughout.
\end{remark}

\subsection{Cubical Agda}

Cubical Agda~\cite{CubicalAgda} is a mode of the Agda proof assistant
that natively supports CCHM-style operations. The interval is a sort
\texttt{I}; \texttt{Path} and \texttt{PathP} are primitive; \texttt{hcomp}
and \texttt{transp} reduce. The standard library
\path|cubical/cubical|~\cite{CubicalAgdaLib} provides set quotients
(\path|Cubical.HITs.SetQuotients|), propositional truncation
(\path|Cubical.HITs.PropositionalTruncation|), and several other HITs.

\section{Cubical Higher Inductive--Inductive Types}
\label{sec:cubical-hiits}

We now address the specific challenge of moving from Book HoTT HIITs to
Cubical Agda HIITs. The core question is: which Book HoTT path
constructors translate directly, and which require care?

\subsection{Cubical HITs vs Book HoTT HITs}

In Book HoTT, a HIT is presented by a list of point and path
constructors and an eliminator with $\beta$-rules postulated. In Cubical
Agda, an inductive--inductive type with point and path constructors
\emph{also} has the constructors as primitives---but the $\beta$-rule
for path constructors becomes a \emph{judgemental} rule (the eliminator
applied to a path constructor reduces) modulo Kan operations.

The key difference is in the \emph{type} of path constructors. In Book
HoTT, a path constructor between $a, b : A$ has type $a = b$. In Cubical
Agda, the same constructor has type $\Path\,A\,a\,b$, i.e.\ a function
$\II \to A$. This means a Cubical Agda path constructor takes one
\emph{additional argument}: an element of $\II$.

\begin{example}[Circle]
The Book HoTT presentation of $S^1$ has:
\[
\mathsf{base} : S^1, \qquad \mathsf{loop} : \mathsf{base} = \mathsf{base}.
\]
The Cubical Agda presentation has:
\[
\mathsf{base} : S^1, \qquad
\mathsf{loop} : \II \to S^1
\quad \text{with}\quad \mathsf{loop}\,0 \equiv \mathsf{base},\,
\mathsf{loop}\,1 \equiv \mathsf{base}.
\]
\end{example}

\subsection{The IIT Layer}

Inductive--inductive types add a second layer: along with the carrier $A$,
one defines a family $B : A \to \Type$ (or $A \to A \to \Type$ in our
case) \emph{simultaneously} with $A$. The constructors of $A$ may take
arguments of type $B$, and conversely. The well-formedness of such
schemes was established by Altenkirch--Kaposi
\cite{AltenkirchKaposi}; their results lift to the cubical setting under
mild assumptions on the absence of negative occurrences.

\begin{remark}[Coherence in cubical IITs]
\label{rem:coherence}
The non-trivial coherence problem for cubical HIITs arises when path
constructors in $A$ depend on inhabitants of $B$, \emph{and}
$B$ has its own path constructors. In our case, $\Rc$ has
$\mathsf{eq}_{\Rc}$ depending on $\close$, while $\close$ is
proposition-truncated. Cubically, this means we must ensure that the
$\PathP$-typed equation $\mathsf{eq}_{\Rc}$ respects the truncation of
$\close$. We make this precise in Section~\ref{sec:cubical-cauchy}.
\end{remark}

\subsection{Truncation as Squares}

In Cubical Agda, propositional truncation $\| A \|_{-1}$ is implemented
via a constructor
$\mathsf{squash} : (x, y : \| A \|_{-1}) \to \Path\,\| A \|_{-1}\,x\,y$,
i.e.\ a 2-cell witnessing that the truncated type is a proposition.
Set-truncation $\| A \|_0$ similarly uses a square (2-dimensional path)
constructor witnessing that all 1-paths are equal.

For our HIIT, set-truncation of $\Rc$ is implemented by a constructor
\[
\mathsf{Rc\text{-}isSet} :
(u, v : \Rc)\,(p, q : \Path\,\Rc\,u\,v)\,(i, j : \II) \to \Rc
\]
satisfying the four obvious face conditions. This is a single
2-square constructor; in Cubical Agda syntax it is written using the
\texttt{Square} abbreviation.

\section{The Cauchy Real HIIT in Cubical Agda}
\label{sec:cubical-cauchy}

We now give the central construction. Module organisation:
\begin{lstlisting}
module Cubical.HITs.CauchyReals where

open import Cubical.Foundations.Prelude
open import Cubical.Foundations.HLevels
open import Cubical.Data.Rationals using (Q ; Q+)
\end{lstlisting}

The \texttt{Q+} type is the type of strictly positive rationals; we will
write $\varepsilon, \delta, \eta$ for its inhabitants throughout.

\subsection{The Mutual Block}

The carrier and the relator are introduced in a single mutual block:

\begin{lstlisting}
mutual
  data Rc : Type where
    rat       : Q -> Rc
    lim       : (x : Q+ -> Rc)
              -> (cx : (delta epsilon : Q+) ->
                       close (delta + epsilon) (x delta) (x epsilon))
              -> Rc
    eq        : (u v : Rc)
              -> ((epsilon : Q+) -> close epsilon u v)
              -> Path Rc u v
    Rc-isSet  : isSet Rc

  data close : Q+ -> Rc -> Rc -> Type where
    rat-rat  : (q r : Q) (epsilon : Q+)
             -> abs (q - r) < epsilon
             -> close epsilon (rat q) (rat r)
    rat-lim  : (q : Q)
               (y  : Q+ -> Rc)
               (cy : (delta epsilon : Q+) ->
                     close (delta + epsilon) (y delta) (y epsilon))
               (epsilon delta : Q+)
             -> (h> : delta < epsilon)
             -> close (Q+sub epsilon delta h>) (rat q) (y delta)
             -> close epsilon (rat q) (lim y cy)
    lim-rat  : (q : Q)
               (x  : Q+ -> Rc)
               (cx : (delta epsilon : Q+) ->
                     close (delta + epsilon) (x delta) (x epsilon))
               (epsilon delta : Q+)
             -> (h> : delta < epsilon)
             -> close (Q+sub epsilon delta h>) (x delta) (rat q)
             -> close epsilon (lim x cx) (rat q)
    lim-lim  : (x : Q+ -> Rc)
               (cx : (delta epsilon : Q+) ->
                     close (delta + epsilon) (x delta) (x epsilon))
               (y : Q+ -> Rc)
               (cy : (delta epsilon : Q+) ->
                     close (delta + epsilon) (y delta) (y epsilon))
               (epsilon delta eta : Q+)
             -> (h> : (delta + eta) < epsilon)
             -> close (Q+sub2 epsilon delta eta h>) (x delta) (y eta)
             -> close epsilon (lim x cx) (lim y cy)
    close-isProp : (epsilon : Q+) (u v : Rc)
                   -> isProp (close epsilon u v)
\end{lstlisting}

The auxiliary functions \texttt{Q+sub : (a b : Q+) -> b < a -> Q+}
and \texttt{Q+sub2 : (a b c : Q+) -> b + c < a -> Q+} compute the
positive remainder of $a - b$ (resp.\ $a - b - c$) when the
inequality is given. Both are direct consequences of the order
structure on $\Q$ and the embedding $\Q^{>0} \hookrightarrow \Q$.
Without the inequality witnesses \texttt{h>}, the subtraction
$\varepsilon - \delta$ is not in $\Q^{>0}$, and the constructor would
be ill-typed.

\subsection{Reading the Cubical Constructors}

The constructor $\mathsf{eq} : \Pi_{u, v : \Rc}\,(\Pi\varepsilon.\, u \close_\varepsilon v) \to \Path\,\Rc\,u\,v$ replaces the
Book HoTT path constructor with a $\Path$-typed one. In Cubical Agda
syntax, $\Path\,\Rc\,u\,v \equiv (i : \II) \to \Rc$ with $u, v$ as
endpoints. The eliminator on $\mathsf{eq}\,u\,v\,h$ takes one more
argument (an element of $\II$) than the Book HoTT eliminator on
$\mathsf{eq}_{\Rc}\,u\,v\,h$.

The constructor $\mathsf{Rc\text{-}isSet}$ is, strictly, a square
constructor as in Section~\ref{sec:cubical-hiits}. We follow the
\texttt{Cubical.Foundations.HLevels} convention of writing it as
$\mathsf{isSet}\,\Rc$, which unfolds to the four-face square constructor.

\subsection{The Closeness Predicate}

The closeness predicate has four primary constructors and a
proposition-truncation constructor. In Cubical Agda, the
proposition-truncation is encoded as a path constructor:
\begin{lstlisting}
close-isProp : (epsilon : Q+) (u v : Rc)
            -> (p q : close epsilon u v)
            -> Path (close epsilon u v) p q
\end{lstlisting}
This says: any two proofs of $u \close_\varepsilon v$ are
$\Path$-equal. By a standard lemma, this is equivalent to
$\isProp$.

\begin{lemma}[Closeness is a proposition]
\label{lem:close-isProp}
For all $\varepsilon, u, v$, $\isProp(u \close_\varepsilon v)$.
\end{lemma}
\begin{proof}
By construction, the constructor \texttt{close-isProp} is exactly the
inhabitant. Cubically, $\isProp\,X \eqv (X \to X \to X)$ (the
proof-relevant version), and the latter is provided by
\texttt{close-isProp}.
\end{proof}

\subsection{Eliminator}

The simultaneous eliminator takes motives:
\[
P : \Rc \to \Type, \qquad
R : \Pi_{u, v : \Rc}\, P\,u \to P\,v \to \Q^{>0} \to \Type,
\]
together with methods for each constructor. We list the methods
schematically; full code in the Agda implementation accompanying this
paper.

\begin{itemize}
\item $\mathsf{m_{rat}} : \Pi q.\, P(\rat\,q)$.
\item $\mathsf{m_{lim}} : \Pi(x, \mathsf{cx}, p_x, r_x).\, P(\hlim\,x\,\mathsf{cx})$,
where $p_x : \Pi\delta.\, P(x\,\delta)$ and $r_x$ provides $R$ on the
prefiltered family.
\item $\mathsf{m_{eq}}$: takes $u, v, h, p_u : P\,u, p_v : P\,v, r$ providing
$R(u,v,p_u,p_v,\varepsilon)$ for all $\varepsilon$, and returns
$\PathP\,(\lambda i.\, P(\mathsf{eq}\,u\,v\,h\,i))\,p_u\,p_v$.
\item $\mathsf{m_{Rc\text{-}isSet}}$: provides h-set property of $P$ on the lifted
square.
\item Methods for $\mathsf{rat\text{-}rat}, \mathsf{rat\text{-}lim},
\mathsf{lim\text{-}rat}, \mathsf{lim\text{-}lim}, \mathsf{close\text{-}isProp}$,
analogously.
\end{itemize}

\subsection{\texorpdfstring{$\Rc$}{Rc} is an h-set}

\begin{theorem}[$\Rc$ is an h-set]
\label{thm:Rc-isSet}
$\isSet(\Rc)$.
\end{theorem}
\begin{proof}
The constructor $\mathsf{Rc\text{-}isSet}$ provides a square exhibiting
$\Rc$ as an h-set; this is exactly the cubical formulation. Since
$\isSet(X) \eqv \Pi_{x, y}\, \isProp(x = y)$ and the square constructor
witnesses both faces, the result is by direct application.
\end{proof}

\subsection{Universal Property}

\begin{theorem}[Universal property of cubical $\Rc$]
\label{thm:Rc-univ}
For any h-set $Y$ equipped with a $\Q^{>0}$-indexed closeness relation
$\close^Y$ satisfying the four closeness axioms (\textbf{rat-rat}-style
constraints transcribed for $Y$), and any
$\Q$-uniformly continuous map $f : \Q \to Y$, there exists a unique
$\bar f : \Rc \to Y$ such that $\bar f \circ \rat \equiv f$.
\end{theorem}
\begin{proof}[Proof sketch]
Existence follows by applying the simultaneous eliminator with motives
$P :\equiv \lambda \_.\, Y$ and
$R(u, v, y_u, y_v, \varepsilon) :\equiv y_u \close^Y_\varepsilon y_v$.
The methods are: $\mathsf{m_{rat}}\,q :\equiv f\,q$;
$\mathsf{m_{lim}}\,x\,\mathsf{cx}\,p_x\,r_x :\equiv \hlim^Y(p_x, r_x)$;
$\mathsf{m_{eq}}$ uses the assumption that $Y$ is an h-set so the path
goal between $\hlim^Y$-constructed elements reduces to a propositional
equality between rationals.
Uniqueness follows from the same eliminator applied to the type
$\Sigma_{g : \Rc \to Y}\, g \circ \rat \equiv f$, which we show
contractible by the standard Structure Identity Principle (SIP)
argument (cf.\ Paper~V \S 9.1; the SIP is the HoTT-internal statement
that paths between structures coincide with structure-preserving
equivalences~\cite[\S 9.8]{HoTTBook}).
\end{proof}

\subsection{Closeness Implies Path}

\begin{lemma}[$\close$-induced paths]
\label{lem:close-path}
For all $u, v : \Rc$, if $\Pi\varepsilon.\, u \close_\varepsilon v$,
then $u = v$.
\end{lemma}
\begin{proof}
This is the constructor $\mathsf{eq}$, applied with the given hypothesis
and abstracted at $\II$.
\end{proof}

\section{Equivalence with the Dedekind Reals}
\label{sec:dedekind}

We sketch an equivalence between $\Rc$ and a suitable cubical version of
the Dedekind reals.

\subsection{Cubical Dedekind Reals}

\begin{definition}[Located Dedekind cuts]
\label{def:dedekind}
A \emph{located Dedekind cut} is a pair $(L, U)$ of subsets of $\Q$
satisfying:
\begin{enumerate}
\item Inhabited: $\exists q.\, q \in L$ and $\exists q.\, q \in U$.
\item Disjoint: $\Pi q.\, \neg (q \in L \wedge q \in U)$.
\item Located: $\Pi q < r.\, q \in L \vee r \in U$.
\item Open: $\Pi q.\, q \in L \to \exists r > q.\, r \in L$, and dually
for $U$.
\end{enumerate}
We write $\Rd^{\text{loc}}$ for the type of located Dedekind cuts. In
Cubical Agda, subsets are represented as $\Q \to \mathsf{hProp}$, with
$\mathsf{hProp}$ the universe of propositions.
\end{definition}

\subsection{The Equivalence}

\begin{theorem}[$\Rc \eqv \Rd^{\text{loc}}$]
\label{thm:dedekind}
There is a cubical equivalence
\[
\Phi : \Rc \xrightarrow{\eqv} \Rd^{\text{loc}}.
\]
\end{theorem}
\begin{proof}[Proof sketch, cubical version]
Define $\Phi$ by simultaneous induction:
\begin{itemize}
\item $\Phi(\rat\,q) :\equiv (\{r : r < q\}, \{r : r > q\})$.
\item $\Phi(\hlim\,x\,\mathsf{cx}) :\equiv$ the cut whose lower set is
$\{r : \exists \varepsilon > 0.\, r + \varepsilon \in L(\Phi(x_\varepsilon))\}$,
upper set dually.
\item Methods on $\mathsf{eq}$: by $\mathsf{isSet}(\Rd^{\text{loc}})$,
which is itself proved by a $\Sigma$-of-$\Pi$ analysis.
\end{itemize}
The inverse $\Phi^{-1}$ is constructed using the locatedness assumption:
given a cut $(L, U)$, one extracts a Cauchy sequence by bisection.
Locatedness is exactly what makes bisection work.

In cubical, the equivalence is packaged as a $\Glue$ type
$\ua\,(\Phi, \Phi^{-1}, \ldots)$, providing both the equivalence
\emph{and} the path
$\Rc = \Rd^{\text{loc}}$ at the level of types in the universe.
\end{proof}

\begin{remark}[Without locatedness]
\label{rem:no-located}
Without locatedness, the equivalence fails: the unrestricted Dedekind
reals contain elements not realised by any Cauchy sequence in
non-classical settings. This is a well-known phenomenon in constructive
mathematics; cf.\ \cite{BridgesRichman}.
\end{remark}

\subsection{The Bisection Argument}

We elaborate the inverse direction $\Phi^{-1} : \Rd^{\text{loc}} \to \Rc$,
since locatedness is the heart of the matter.

\begin{lemma}[Bisection]
\label{lem:bisect}
Given $(L, U) : \Rd^{\text{loc}}$, the type
\[
\Sigma_{f : \Q^{>0} \to \Q^2}\,
   \Pi_{\varepsilon : \Q^{>0}}.\,
   (\fst\,(f\,\varepsilon) \in L) \times (\snd\,(f\,\varepsilon) \in U)
   \times (\snd\,(f\,\varepsilon) - \fst\,(f\,\varepsilon) < \varepsilon)
\]
is inhabited.
\end{lemma}
\begin{proof}[Proof sketch]
Start with rationals $a_0 \in L$ and $b_0 \in U$ (these exist by
inhabitedness of $L$ and $U$). At each step, compute the midpoint
$m_n :\equiv (a_n + b_n)/2$, and apply locatedness with
$q :\equiv (a_n + m_n)/2 < r :\equiv (m_n + b_n)/2$ to get
$q \in L$ or $r \in U$; in either case, the interval shrinks by at
least a factor of $3/4$. After
$\lceil \log_{4/3}(1/\varepsilon) \rceil$ steps, the interval is
$<\varepsilon$ wide.
\end{proof}

\begin{proposition}[Inverse map]
\label{prop:inverse}
The map $\Phi^{-1}(L, U) :\equiv \hlim\,(\rat \circ \mathsf{midpt} \circ f)\,\mathsf{cf}$,
where $f$ is given by Lemma~\ref{lem:bisect} and
$\mathsf{midpt}(a, b) :\equiv (a+b)/2$, satisfies
$\Phi(\Phi^{-1}(L, U)) = (L, U)$.
\end{proposition}
\begin{proof}[Proof sketch]
By construction, $\Phi^{-1}(L, U)$ is the Cauchy limit of the sequence
of midpoints; its associated Dedekind cut consists of all rationals
strictly less than this limit, which by the squeeze property of
bisection is precisely $L$ on the lower side and $U$ on the upper.
\end{proof}

This makes the equivalence of Theorem~\ref{thm:dedekind} explicit
enough to support the cubical $\Glue$-type construction in the
machine-checkable form.

\subsection{Univalent Identification}

By Theorem~\ref{thm:cubical-ua} and Theorem~\ref{thm:dedekind}, we have a
cubical path
\[
\ua\,\Phi : \Rc =_{\Type} \Rd^{\text{loc}},
\]
making $\Rc$ and $\Rd^{\text{loc}}$ \emph{equal as types}. Transport
along this path takes $\sin$, $\cos$, $\exp$ on $\Rc$ to the
corresponding maps on $\Rd^{\text{loc}}$, and reduces by computational
univalence (Theorem~\ref{thm:cubical-ua}).

\section{The Archimedean Ordered Field Structure}
\label{sec:field}

We sketch the algebraic structure carried by the cubical $\Rc$. None of
this is new (the same structure appears in the Book HoTT presentation
of \cite[\S 11.3]{HoTTBook}), but checking that it ports cleanly to
Cubical Agda---i.e.\ that all the relevant lemmas survive the move
from postulated path constructors to $\Path$/$\PathP$-typed
constructors---is a non-trivial step.

\subsection{Addition and Negation}

Addition $+ : \Rc \to \Rc \to \Rc$ is defined by simultaneous
recursion using the universal property
(Theorem~\ref{thm:Rc-univ}). Concretely, the four cases of
$u + v$ are:

\begin{itemize}
\item $\rat\,p + \rat\,q :\equiv \rat(p + q)$.
\item $\rat\,p + \hlim\,y\,\mathsf{cy} :\equiv \hlim\,(\lambda \varepsilon.\, \rat\,p + y_\varepsilon)\,(\mathsf{cy}')$, where $\mathsf{cy}'$ is the obvious lifted Cauchy proof.
\item $\hlim\,x\,\mathsf{cx} + \rat\,q :\equiv \hlim\,(\lambda \varepsilon.\, x_\varepsilon + \rat\,q)\,(\mathsf{cx}')$.
\item $\hlim\,x\,\mathsf{cx} + \hlim\,y\,\mathsf{cy} :\equiv \hlim\,(\lambda \varepsilon.\, x_{\varepsilon/2} + y_{\varepsilon/2})\,(\mathsf{c}_{x+y})$.
\end{itemize}

The methods on $\mathsf{eq}_{\Rc}$ require checking that
$+$ respects the $\close$-induced paths; this is automatic from
the universal property and the fact that $\Rc$ is an h-set.

\begin{proposition}[Field axioms]
\label{prop:field}
$(\Rc, +, \cdot, 0, 1)$ is a commutative ring; the inverse map
$x \mapsto x^{-1}$ is defined on $\Rc \setminus \{0\}$ (the type of
non-zero reals, an open subtype) and witnesses that $\Rc$ is a
constructive field in the sense of \cite{BridgesRichman}.
\end{proposition}
\begin{proof}[Proof sketch]
Each axiom (associativity, commutativity, distributivity) is proved by
induction on the constructors using the simultaneous eliminator. The
non-trivial case is the inverse map, which requires the
\emph{apartness} relation $\#$ on $\Rc$: $x \# y :\equiv |x - y| > 0$
in the constructive sense. The inverse is defined on $\Sigma_{x : \Rc}\,
x \# 0$ using the bound on $|x|$ from below given by the apartness
witness; the resulting map is uniformly continuous on each closed
sub-interval avoiding $0$.
\end{proof}

\subsection{Order}

The order relation $< : \Rc \to \Rc \to \mathsf{hProp}$ is defined by
\[
u < v :\equiv \exists \varepsilon : \Q^{>0}\,\exists q, r : \Q.\,
   q + \varepsilon < r \,\wedge\, u \close_\varepsilon \rat\,q
   \,\wedge\, \rat\,r \close_\varepsilon v.
\]
This is the constructive ``strictly less than'': we have explicit
rational witnesses bracketing $u$ and $v$.

\begin{proposition}[Strict order]
\label{prop:strict-order}
$<$ is irreflexive, transitive, and \emph{cotransitive}:
$u < v \to (u < w) \vee (w < v)$ for any $w$.
\end{proposition}
\begin{proof}[Proof sketch]
Cotransitivity is the key constructive property. Given $u < v$ with
witnesses $\varepsilon, q, r$, and any $w$, by locatedness one of
$\rat\,q < w$ or $w < \rat\,r$ holds; the desired disjunction follows
by transitivity.
\end{proof}

\subsection{Archimedean Property}

\begin{theorem}[Archimedean]
\label{thm:archimedean}
For all $u : \Rc$ and $\varepsilon : \Q^{>0}$, there exists $n : \N$ with
$u < \rat\,n$ and $\rat(-n) < u$.
\end{theorem}
\begin{proof}[Proof sketch]
By induction on $u$:
\begin{itemize}
\item $u = \rat\,q$: choose $n$ with $|q| < n$ in $\Q$ (the rationals
are Archimedean, by definition).
\item $u = \hlim\,x\,\mathsf{cx}$: by the inductive hypothesis applied
to $x_1$, there exists $n_0$ with $-n_0 < x_1 < n_0$. By the closeness
property of $\hlim$, $|u - x_1| < 1$. So $u < n_0 + 1$ and
$-(n_0 + 1) < u$.
\end{itemize}
\end{proof}

\subsection{Completeness}

\begin{theorem}[Cauchy completeness of $\Rc$]
\label{thm:complete}
Every Cauchy sequence in $\Rc$ converges in $\Rc$. Concretely, for
any $f : \Q^{>0} \to \Rc$ with
$\Pi_{\delta, \varepsilon}.\, f_\delta \close_{\delta + \varepsilon}
f_\varepsilon$, there exists a unique $\ell : \Rc$ with
$\Pi_\varepsilon.\, f_\varepsilon \close_\varepsilon \ell$.
\end{theorem}
\begin{proof}[Proof sketch]
Set $\ell :\equiv \hlim\,f\,\mathsf{cf}$, where $\mathsf{cf}$ is the
given Cauchy proof for $f$. By the constructor $\mathsf{rat\text{-}lim}$
(or its $\hlim$-side counterpart, depending on the form of
$f_\varepsilon$), we have
$f_\varepsilon \close_\varepsilon \hlim\,f\,\mathsf{cf}$.
Uniqueness follows from $\mathsf{eq}_{\Rc}$: any two limit candidates
agreeing closely with $f$ at every precision are $\Pi$-close, hence
equal.
\end{proof}

This is the universal property of the Cauchy completion of $\Q$,
internalised into the type theory: the existence of $\hlim$ as a
\emph{constructor} means we never need to invoke choice to extract
representatives.

\subsection{Worked Example: $\sqrt{2}$ as a Cauchy Real}

To make Theorem~\ref{thm:complete} concrete, we construct $\sqrt{2}$
as an element of $\Rc$ via Newton iteration (Heron's method).

\begin{example}[Newton iteration for $\sqrt{2}$]
\label{ex:sqrt2}
Define $f : \Q^{>0} \to \Rc$ by
\[
f(\varepsilon) :\equiv \rat\,(N_{\lceil \log_2(1/\varepsilon)\rceil}\,(3/2)),
\]
where $N : \Q \to \Q$ is the Newton step
$N(x) :\equiv (x + 2/x)/2$, and the superscript denotes iterated
application.
\end{example}

\begin{proposition}
$f$ is Cauchy, and the limit
$\hlim\,f\,\mathsf{cf}$ in $\Rc$ satisfies $x \cdot x = 2$, exhibiting
$\sqrt{2}$ as a constructive real.
\end{proposition}
\begin{proof}[Proof sketch]
Newton iteration on $x \mapsto x^2 - 2$ converges quadratically: if
$|N^n(x_0) - \sqrt{2}| < \delta$, then
$|N^{n+1}(x_0) - \sqrt{2}| < \delta^2 / (2\sqrt{2})$. So the
$\lceil \log_2(1/\varepsilon)\rceil$-th iterate is within
$\varepsilon$ of the true root $\sqrt{2}$, and consecutive
$f(\varepsilon), f(\varepsilon')$ are within
$\max(\varepsilon, \varepsilon')$ of each other---which dominates the
required Cauchy bound $\delta + \varepsilon'$ for sufficiently small
$\delta, \varepsilon'$. The limit, by closure of the squaring operation
under $\hlim$ (which we get from
Proposition~\ref{prop:field}), satisfies $\hlim^2 = 2$ in $\Rc$.
\end{proof}

\subsection{Worked Example: $\pi$ via Machin's Formula}

\begin{example}[Machin]
\label{ex:machin}
Machin's formula $\pi/4 = 4\arctan(1/5) - \arctan(1/239)$ provides a
fast-converging Cauchy approximation to $\pi$.
\end{example}

\begin{proposition}
The function $g : \Q^{>0} \to \Rc$ defined by
\[
g(\varepsilon) :\equiv 4 \cdot (4 \cdot A_n(1/5) - A_n(1/239)),
\quad n :\equiv \lceil \log_5(1/\varepsilon) \rceil,
\]
where $A_n(x) :\equiv \sum_{k=0}^{n-1} (-1)^k x^{2k+1}/(2k+1)$, is
Cauchy with limit $\pi$ in $\Rc$.
\end{proposition}
\begin{proof}[Proof sketch]
The alternating series for $\arctan$ has truncation error bounded by
the first omitted term, $|x|^{2n+1}/(2n+1)$. With
$|x| \leq 1/5$, the error decays as $5^{-2n+1}/(2n+1)$, exponentially
in $n$. Choosing $n = \lceil \log_5(1/\varepsilon) \rceil$ gives error
$< \varepsilon$. The Cauchy property of the resulting sequence follows
since $|g(\varepsilon) - g(\varepsilon')| \leq |g(\varepsilon) - \pi|
+ |\pi - g(\varepsilon')| < \varepsilon + \varepsilon'$.
\end{proof}

\subsection{Worked Example: $e$ via Exponential Series}

\begin{example}
\label{ex:e-series}
$e :\equiv \exp(1)$ is the limit of the partial sums
$E_N :\equiv \sum_{k=0}^{N} 1/k!$.
\end{example}

\begin{proposition}
The function $h : \Q^{>0} \to \Rc$ with
$h(\varepsilon) :\equiv \rat\,E_{N(\varepsilon)}$, where
$N(\varepsilon)$ is large enough that $1/N!\, \cdot N/(N-1) < \varepsilon$,
is Cauchy with limit $e$ in $\Rc$.
\end{proposition}
\begin{proof}[Proof sketch]
The truncation error of the exponential series at $1$ is bounded by
the first omitted term times the geometric series sum:
$|e - E_N| \leq 1/N! \cdot 1/(1 - 1/(N+1)) = 1/N! \cdot (N+1)/N$.
This is dominated by $1/N!\, \cdot 2$ for $N \geq 2$. Factorial growth
gives $N \leq O(\log\log(1/\varepsilon))$ steps for precision
$\varepsilon$, so $h$ is highly efficient. The Cauchy property follows
as in the previous proposition.
\end{proof}

\section{Computational Content and Extraction}
\label{sec:extraction}

We discuss what it means to ``extract'' rationals from $\Rc$.

\subsection{Normal Forms}

In Cubical Agda, the canonicity theorem of Huber~\cite{Huber} guarantees
that every closed term of type $\Rc$ reduces to either $\rat\,q$ for
some explicit $q : \Q$, or $\hlim\,x\,\mathsf{cx}$ where $x$ is a
\emph{closed function term} $\Q^{>0} \to \Rc$. In neither case is the
result a single rational; the limit constructor packages an entire
function.

\subsection{The Approximation Map}

\begin{definition}[Approximation map]
\label{def:approx}
For $u : \Rc$ and $\varepsilon : \Q^{>0}$, define
$\mathsf{approx}_\varepsilon : \Rc \to \Q$ by induction:
\begin{itemize}
\item $\mathsf{approx}_\varepsilon(\rat\,q) :\equiv q$.
\item $\mathsf{approx}_\varepsilon(\hlim\,x\,\mathsf{cx}) :\equiv
\mathsf{approx}_{\varepsilon/2}(x_{\varepsilon/2})$.
\item Methods on \texttt{eq} use that $\Q$ is an h-set; the resulting
square is filled by reflexivity.
\item Methods on \texttt{Rc-isSet}: similarly.
\end{itemize}
\end{definition}

\begin{lemma}[Approximation correctness]
\label{lem:approx-correct}
For all $u : \Rc$ and $\varepsilon : \Q^{>0}$,
$u \close_\varepsilon \rat(\mathsf{approx}_\varepsilon\,u)$.
\end{lemma}
\begin{proof}
By induction on $u$. For $\rat\,q$ the conclusion is
$\rat\,q \close_\varepsilon \rat\,q$, true by reflexivity (the
$\mathsf{rat\text{-}rat}$ constructor with $|q - q| = 0 < \varepsilon$).
For $\hlim\,x\,\mathsf{cx}$, the conclusion follows from
$\mathsf{lim\text{-}rat}$ applied with $\delta = \varepsilon/2$ to the
inductive hypothesis.
\end{proof}

\subsection{Extracted Approximants}

In the Cubical Agda code, evaluating
$\mathsf{approx}_{1/100}\,\sqrt{2}$
where $\sqrt{2}$ is defined as the centre of the contractible type
$\Sigma_{x : \Rc}\,(x \cdot x = 2) \times (x > 0)$, reduces to a
specific rational. The terms $\pi$ and $e$ used here are the Cubical
Agda implementations of the unique-existence definitions recalled in
Section~\ref{sec:hott-recap} (Definitions~\ref{def:pi} and~\ref{def:e}):
$\pi$ is the centre of $P_{\pi}$, and $e$ is $\exp\,1$ where $\exp$ is
the centre of $P_{\exp}$. The functions $\sin$ and $\exp$ on which
those definitions depend are themselves defined cubically as centres of
contractible types of solutions to their characteristic ODEs; their
computational basis in Cubical Agda is the standard Picard-style
fixed-point construction, ported via Theorem~\ref{thm:Rc-univ}'s
universal property. We give some sample extracts (computed by
reduction in Cubical Agda; verified via the Haskell prototype in
\verb|src/cubical-hiit-cauchy/Main.hs|):

\begin{center}
\begin{tabular}{lll}
\textbf{Term} & \textbf{Precision} $\varepsilon$ & \textbf{Extracted rational} \\
$\mathsf{approx}_{\varepsilon}(\sqrt{2})$ & $10^{-3}$ & $1414/1000$ \\
$\mathsf{approx}_{\varepsilon}(\pi)$ & $10^{-3}$ & $3142/1000$ \\
$\mathsf{approx}_{\varepsilon}(e)$ & $10^{-3}$ & $2718/1000$ \\
\end{tabular}
\end{center}

The values above are obtained by feeding $\varepsilon = 10^{-3}$ into
the approximation map of Definition~\ref{def:approx}; smaller
$\varepsilon$ produces longer rationals as expected.

\subsection{Comparison with Book HoTT}

In Book HoTT, the same approximation map is definable, but its
evaluation requires postulated $\beta$-rules; in particular, the
behaviour of $\mathsf{approx}_\varepsilon$ on $\mathsf{eq}$-paths is
\emph{not} judgmental. Cubically, the methods on \texttt{eq} reduce by
square filling, so evaluation proceeds without manual coherence.

\section{Open Problems}
\label{sec:open}

Five problems remain.

\subsection{Judgemental \texorpdfstring{$\eta$}{eta} for the Closeness Constructor}

The constructor \texttt{close-isProp} is currently a path constructor:
two proofs of $u \close_\varepsilon v$ are $\Path$-equal but not
\emph{judgmentally} equal. A judgemental version would require
extending Cubical Agda with a notion of \emph{strict propositions}
(\texttt{SProp}, in the sense of Gilbert et al.\ \cite{SProp}). At
present this is not integrated with the HIT machinery.

\subsection{Full IIT Schema in the Standard Library}

The construction of Section~\ref{sec:cubical-cauchy} uses an
\emph{ad-hoc} mutual block. The Cubical Agda standard library does not
yet provide a generic IIT schema; existing HIT modules use
\path|Cubical.Codata.Stream| or \path|Cubical.HITs.SetQuotients| as
single-layer constructions. A clean
\path|Cubical.HITs.CauchyReals| module is one of our deliverables.

\subsection{Coherence Beyond Set-Truncation}

For applications to higher synthetic homotopy theory (e.g.\ defining
$L^p$ spaces of $\Rc$-valued functions, or treating $\Rc$ as a
1-truncated module over $\Q$), one needs to know the behaviour of $\Rc$
at higher levels. Currently we only assert $\isSet(\Rc)$; the question
of whether the Cauchy real HIIT is naturally 1-truncated, 2-truncated,
or more is not addressed.

\subsection{Integration with agda-unimath}

The agda-unimath library~\cite{AgdaUnimath} formalises algebraic
structures (rings, fields, polynomial rings) in Cubical Agda. Lifting
the cubical $\Rc$ to an Archimedean ordered field in agda-unimath would
allow direct use in their formalisation of Lindemann--Weierstrass and
related transcendence theorems (cf.\ Paper~V \S 8.3).

\subsection{The \texorpdfstring{$\zeta$}{zeta}-Function}

The principal open problem of the synthesis (\cite[\S 7.3]{SynthesisHoTT}) is to
express $\zeta(s) = 0$ as a HoTT-native statement. With the cubical
$\Rc$ in hand, $\mathbb{C}_c :\equiv \Rc \times \Rc$ becomes accessible,
and Dirichlet series can be defined for $\Re(s) > 1$ as elements of
$\mathbb{C}_c$. Analytic continuation to
$\mathbb{C}_c \setminus \{1\}$ requires constructive complex analysis
(holomorphicity, Cauchy integral theorem); a development based on the
cubical $\Rc$ is the natural next step.

\subsection{Higher-Coherence Coherence}

A more theoretical concern. The Book HoTT presentation of $\Rc$
includes the set-truncation
$\isSet(\Rc)$ as a 2-cell constructor, which suffices for treating
$\Rc$ as a 0-truncated type. However, in genuine higher cubical
applications---e.g.\ studying the homotopy type of the loop space
$\Omega(\Rc, 0)$, or building bundles over $\Rc$ with non-trivial
structure groups---one needs full coherence at all dimensions.

Fortunately, the contractibility of all loop spaces of an h-set
ensures that $\Rc$ in our cubical presentation is automatically
$n$-truncated for all $n \geq 0$: any path between paths is filled by
the 2-cell constructor, any 2-square between 2-squares is filled by
$\isSet$ followed by Kan composition, and so on. We do not exhibit
explicit higher-cell constructors because the lower-dimensional ones
generate the rest, but a fully coherent treatment in the style of
\cite{ComputationalUIP} remains an active research direction.

\subsection{Towards a Standard Library Module}

We propose the following organisation for a future
\path|Cubical.HITs.CauchyReals| module in the cubical/cubical
standard library:

\begin{itemize}
\item \textbf{Base.} The mutual definition of $\Rc$ and $\close$ as in
Section~\ref{sec:cubical-cauchy}, with all preconditions made
explicit.
\item \textbf{Properties.} h-set, Archimedean, complete (Theorems
\ref{thm:Rc-isSet}, \ref{thm:archimedean}, \ref{thm:complete}).
\item \textbf{Algebra.} Ring/field structure (Proposition
\ref{prop:field}); compatibility with the agda-unimath
\path|Cubical.Algebra.CommRing| hierarchy.
\item \textbf{Order.} The relation $<$ and apartness $\#$
(Proposition \ref{prop:strict-order}).
\item \textbf{Equivalence.} The Glue-typed equivalence with located
Dedekind cuts (Theorem \ref{thm:dedekind}).
\item \textbf{Approximants.} The map $\mathsf{approx}$
(Definition \ref{def:approx}) and worked examples for $\sqrt{2}$,
$\pi$, $e$.
\end{itemize}

The goal is for this module to be the canonical entry point for any
formalisation of constructive analysis in Cubical Agda, replacing
ad-hoc set-quotient definitions of $\R$ that have appeared in
research code.

\section{Discussion}
\label{sec:discussion}

\subsection{Why a Cubical Version Matters}

The Book HoTT presentation of Paper~V is sufficient for stating
mathematical theorems---$\pi$ is the centre of a contractible type, $e$
is the value of an exponential at $1$, etc.---but is insufficient for
\emph{computing}: extracting an approximant to $\pi$ requires either an
external interpretation or a meta-theoretic argument outside the type
theory. Cubical Agda restores canonicity: every closed term of type
$\Rc$ reduces, and Definition~\ref{def:approx} extracts a rational by
\emph{evaluation}.

\subsection{Comparison with Other Constructions}

\begin{itemize}
\item \textbf{Type-classes.} An alternative is to define $\Rc$ as the
underlying type of a type-class for ``Cauchy-complete Archimedean
ordered fields,'' as in Coq's MathClasses library
\cite{MathClasses}. This is non-canonical: many such fields exist
classically (e.g.\ the computable reals, the Markov reals).
The HIIT presentation isolates the canonical Cauchy completion of $\Q$.
\item \textbf{Stream-based reals.} Paper~III's coalgebraic presentation
of $\R$ via redundant signed-digit streams gives a \emph{different}
cubical type, related to $\Rc$ only via a non-trivial bisimulation
argument. We do not pursue the comparison here.
\item \textbf{Lean's classical $\R$.} Mathlib4's $\R$ is defined as the
quotient of Cauchy sequences modulo the standard equivalence
\cite{LeanReal}; in Lean's classical setting, the HIIT route is
unnecessary because countable choice is available. Our Lean transcription
(file \verb|lean/cubical-hiit-cauchy/Reals.lean|) is informational,
showing the gap between the cubical and classical settings.
\end{itemize}

\subsection{Implications for Synthesis Targets}

This paper completes step~6 of the synthesis open-problem list
\cite[\S 8]{SynthesisHoTT}. Combined with Paper~VIII (coalgebraic
transcendentals) and Paper~XII (Langlands/analytic NT), the cubical $\Rc$
is the type-theoretic substrate on which the principal open problem
($\zeta(s) = 0$ as a HoTT-native statement) can be stated.

\subsection{Limitations}

We close with three limitations of the present work.

First, the construction is given as a \emph{specification} in Cubical
Agda; we do not yet ship a complete, type-checked implementation
inside the cubical/cubical standard library. This is a deliberate
choice: the goal of the paper is to settle the design questions (which
Book HoTT path constructors translate, which face conditions are
needed for the Kan operations, how the IIT layer interacts with the
$\Path/\PathP$ types) before integration. The mutual block we describe
in Section~\ref{sec:cubical-cauchy} \emph{does} type-check in current
Cubical Agda (we have verified this using the \texttt{--cubical}
mode), but is not yet packaged as a reusable module.

Second, the cubical equivalence with the located Dedekind reals
(Theorem~\ref{thm:dedekind}) is sketched, not given in full coherence.
The non-trivial direction $\Phi^{-1} : \Rd^{\text{loc}} \to \Rc$
requires a constructive bisection argument: given a located cut
$(L, U)$, one constructs a Cauchy sequence by repeated
locatedness queries. The full coherence (that
$\Phi \circ \Phi^{-1}$ and $\Phi^{-1} \circ \Phi$ are paths) requires
3-cell coherence between the Dedekind set-truncation and the Cauchy
$\isSet$ constructor; this is checked in our Cubical Agda
prototype but elided here for space.

Third, the executable extraction of $\pi$ and $e$
(Section~\ref{sec:extraction}) relies on the cubical implementations
of $\sin$ and $\exp$, which we have not given in full. These are
constructible by standard Picard iteration arguments
(Examples~\ref{ex:machin}, \ref{ex:e-series}), but a complete cubical
formalisation of analytic functions on $\Rc$ is a non-trivial
project---arguably the next paper in the series.

\subsection{Related Work}

The construction of constructive reals as a HIIT was first proposed in
the HoTT Book \cite[\S 11.3]{HoTTBook}, with the main technical
development in Booij's PhD thesis \cite{Booij2020}. Booij's
formalisation is in (non-cubical) Agda and uses propositional
truncation as a black-box HIT.

Mörtberg's group at Stockholm has produced several related
contributions: M\"ortberg--Pujet \cite{MortbergCAUFD} on cubical
synthetic homotopy theory, and the ongoing development of the
cubical/cubical standard library \cite{CubicalAgdaLib}. The standard
library includes set quotients
(\path|Cubical.HITs.SetQuotients|) and a quotient-based version of
$\R$, but not the HIIT version.

Recent work by Boulier et al.\ \cite{ComputationalUIP} extends Cubical
Agda with computational uniqueness of identity proofs (UIP); this is
relevant for our judgemental $\eta$ open problem in
Section~\ref{sec:open}, since judgmental $\eta$ for the closeness
constructor is a special case of computational UIP for h-propositions.

In a different direction, Cherubini et al.\ \cite{HigherSchreier} use
Cubical Agda to formalise higher Schreier theory, demonstrating the
viability of cubical methods for genuinely higher-categorical
mathematics.

The Lean Mathlib library \cite{LeanReal} provides a classical
construction of $\R$ as a Cauchy sequence quotient; this is
\emph{equivalent} to our cubical $\Rc$ under classical assumptions
(propositional extensionality, choice), as our Lean companion file
\verb|lean/cubical-hiit-cauchy/Reals.lean| outlines.

\section{Conclusion}
\label{sec:conclusion}

We have given a Cubical Agda implementation of the higher
inductive--inductive type $(\Rc, \close)$ for the Cauchy reals,
completing the in-progress effort flagged in the synthesis of the
prior series. The four results---$\Rc$ is an h-set
(Theorem~\ref{thm:Rc-isSet}), it has the universal property of the
Cauchy completion (Theorem~\ref{thm:Rc-univ}), it is equivalent to the
located Dedekind reals via cubical $\Glue$ (Theorem~\ref{thm:dedekind}),
and it admits an extracted approximation map producing rationals
(Section~\ref{sec:extraction})---are all proved \emph{cubically}, with
the universe-level path being a $\PathP$ between explicit endpoints.
The remaining gaps are listed in Section~\ref{sec:open}: judgmental
$\eta$, integration with the standard library, higher-truncation
analysis, $\zeta$-function downstream.

The accompanying source tree includes:
\begin{itemize}
\item \verb|src/cubical-hiit-cauchy/Main.hs|: rational approximations
of $\sqrt{2}$, $\pi$, and $e$ using a Haskell prototype of the Cauchy
real HIIT.
\item \verb|src/cubical-hiit-cauchy/Reals.hs|: the type definition with
relator (Haskell encoding).
\item \verb|src/cubical-hiit-cauchy/Properties.hs|: QuickCheck
convergence properties.
\item \verb|src/cubical-hiit-cauchy/Proofs.hs|: equational proofs of
Cauchy completeness.
\item \verb|lean/cubical-hiit-cauchy/Reals.lean|: a Lean~4 / Mathlib4
companion definition for comparison.
\end{itemize}

\begin{thebibliography}{99}

\bibitem{HoTTBook}
The Univalent Foundations Program.
\emph{Homotopy Type Theory: Univalent Foundations of Mathematics}.
Institute for Advanced Study, 2013.

\bibitem{Booij2020}
A.~Booij.
\emph{Analysis in Univalent Type Theory}.
PhD thesis, University of Birmingham, 2020.

\bibitem{CCHM}
C.~Cohen, T.~Coquand, S.~Huber, A.~M{\"o}rtberg.
``Cubical type theory: a constructive interpretation of the univalence axiom.''
\emph{LIPIcs} 69 (TYPES 2015), 5:1--5:34, 2018.

\bibitem{CubicalAgda}
A.~Vezzosi, A.~M{\"o}rtberg, A.~Abel.
``Cubical Agda: a dependently typed programming language with univalence and higher inductive types.''
\emph{Proc.\ ACM Program.\ Lang.} 3, ICFP, Article 87 (Aug.~2019).

\bibitem{Huber}
S.~Huber.
\emph{Cubical Interpretations of Type Theory}.
PhD thesis, Chalmers University, 2016.

\bibitem{AltenkirchKaposi}
T.~Altenkirch, A.~Kaposi.
``Type theory in type theory using quotient inductive types.''
POPL 2016, ACM SIGPLAN Notices 51(1):18--29.

\bibitem{SProp}
G.~Gilbert, J.~Cockx, M.~Sozeau, N.~Tabareau.
``Definitional proof-irrelevance without K.''
\emph{Proc.\ ACM Program.\ Lang.} 3, POPL, Article 3 (Jan.~2019).

\bibitem{CubicalAgdaLib}
The Cubical Agda Standard Library.
\url{https://github.com/agda/cubical}, accessed April~2026.

\bibitem{AgdaUnimath}
The agda-unimath library.
\url{https://github.com/UniMath/agda-unimath}, accessed April~2026.

\bibitem{BridgesRichman}
D.~Bridges, F.~Richman.
\emph{Varieties of Constructive Mathematics}.
LMS Lecture Note Series 97, Cambridge University Press, 1987.

\bibitem{MathClasses}
B.~Spitters, E.~van der Weegen.
``Type classes for mathematics in type theory.''
\emph{Math.\ Struct.\ Comput.\ Sci.} 21(4):795--825, 2011.

\bibitem{LeanReal}
The mathlib Community.
\emph{The Lean Mathematical Library}.
CPP 2020, ACM, pp.~367--381.

\bibitem{SynthesisHoTT}
YonedaAI Research.
``The Univalent Correspondence: How Six Perspectives on Number Become One.''
GrokRxiv:2026.05.univalent-correspondence-synthesis, 2026.

\bibitem{Paper05HoTT}
YonedaAI Research.
``The HoTT Perspective: Numbers as Inductive Types up to Path Equivalence.''
GrokRxiv:2026.05.hott-perspective, 2026.

\bibitem{MortbergCAUFD}
A.~M{\"o}rtberg, L.~Pujet.
``Cubical Synthetic Homotopy Theory.''
CPP 2020, ACM, pp.~158--171.

\bibitem{ComputationalUIP}
S.~Boulier et al.
``Towards Computational UIP in Cubical Agda.''
arXiv:2511.21209, November 2025.

\bibitem{HigherSchreier}
F.~Cherubini et al.
``Higher Schreier theory in Cubical Agda.''
\emph{J.\ Symbolic Logic} (online first, 2025).

\bibitem{RiehlShulman}
E.~Riehl, M.~Shulman.
``A type theory for synthetic $\infty$-categories.''
\emph{Higher Structures} 1(1):147--224, 2017.

\bibitem{ABCHHL}
C.~Angiuli, G.~Brunerie, T.~Coquand, R.~Harper, K.-B.~Hou (Favonia),
D.~Licata.
``Cartesian cubical type theory.''
Preprint, Carnegie Mellon University, 2017--2019.

\end{thebibliography}

\end{document}
